%%%%%%%%%%%%%%%%%%%%%%%%%%%%%%%%%%%%%%%%%%%%%%%%%%%%%%%%%%%%%%%%%%%%%%
% Lớp 11 | Đạo hàm | Các quy tắc đạo hàm
% Dạng 1D7?2-4: Tiếp tuyến biết trước hệ  số góc
% Số câu: 57  (sinh tự động bởi tach_cau.py; từ nay thầy sửa trực tiếp file này)
%%%%%%%%%%%%%%%%%%%%%%%%%%%%%%%%%%%%%%%%%%%%%%%%%%%%%%%%%%%%%%%%%%%%%%

%%%%%%%%%%%%%%%% TRẮC NGHIỆM 4 PHƯƠNG ÁN %%%%%%%%%%%%%%%%

% nguồn: Ok_TracNghiemLop11_Lan1
\begin{ex}%[1D7N2-4]%[ma:a161706f]
	Cho hàm số $y=x^3-2x+1$ có đồ thị $(C)$. Hệ số góc $k$ của tiếp tuyến với $(C)$ tại  điểm có hoàng độ bằng $1$ bằng
	 \choice
		 {$k=-5$}
		 {$k=10$}
		 {$k=25$}
		{\True  $k=1$}
	 \loigiai{
		Ta  có $y'=3x^2-2$.\\
	Hệ số góc $k$ của tiếp  tuyến với $(C)$ tại điểm có hoàng độ bằng $1$ bằng $k=y'(1)=1$.}
\end{ex}

% nguồn: Ok_TracNghiemLop11_Lan1
\begin{ex}%[1D7N2-4]%[ma:92352fbc]
	Tiếp tuyến của đồ thị hàm số $y=x^3+3x$ tại điểm $M(-1;-4)$ có hệ số góc $k$ là 
	 \choice
		{$k=4$} 
		 {$k=3$}
		 {$k=0$}
		 {\True $k=6$}
	\loigiai{ 
		$y=x^3+3x\Rightarrow y'=3x^2+3\Rightarrow y'(-1)=6$. Vậy tiếp tuyến của đồ thị hàm số tại $M$  có hệ số góc $k=y'(-1)=6$.
	} 
\end{ex}

% nguồn: Ok_Toan11_TracNghiem_Lan2
\begin{ex}%[1D7H2-4]%[ma:1bda4cbf]
	Viết phương trình tiếp tuyến của đồ thị hàm số  $y=f(x)=\sqrt{2x+1}$, biết rằng tiếp tuyến đó song song với đường thẳng $x-3y+6=0$.
	 \choice
		 {$y=-\dfrac{1}{3}x-1$}
		{$y=-\dfrac{1}{3}x+1$} 
		{$y=-\dfrac{1}{3}x-\dfrac{5}{3}$} 
		{\True $y=\dfrac{1}{3}x+\dfrac{5}{3}$} 
	 \loigiai{
		Gọi $M(x_0;y_0)$  là tiếp điểm.
		\begin{align*}
			y=\sqrt{2x+1}\Rightarrow y'=f '(x)=\dfrac{1}{\sqrt{2x+1}}. 
		\end{align*}
		Ta có $x-3y+6=0\Leftrightarrow y=\dfrac{1}{3}x+2$. Tiếp tuyến có hệ số góc bằng  $\dfrac{1}{3}$ khi và chỉ khi
		\begin{align*}
			 f'(x_0)=\dfrac{1}{\sqrt{2x_0+1}}=\dfrac{1}{3}\Leftrightarrow\sqrt{2x_0+1}=3\Leftrightarrow 2x_0+1=9\Leftrightarrow x_0=4.
		\end{align*}
		 Với $x_0=4$, ta được $y_0=3$.\\
		Phương trình tiếp  tuyến là
		\begin{align*}
			 y-3=\dfrac{1}{3}(x-4)\Leftrightarrow y=\dfrac{1}{3}x+\dfrac{5}{3}.
		\end{align*}
	 }
\end{ex}

% nguồn: Ok_Toan11_TracNghiem_Lan2
\begin{ex}%[1D7H2-4]%[ma:10f65b43]
	Cho hàm số $y=\dfrac{1}{3}x^3-2x^2+x+2$ có đồ  thị $(C)$. Phương trình các tiếp tuyến với đồ thị $(C)$ biết tiếp tuyến song song với đường thẳng $d\colon y=-2x+\dfrac{10}{3}$ là
	\choice 
		 {\True $y=-2x+2$}
		{$y=-2x-2$} 
		{$y=-2x+10$,  $y=-2x-\dfrac{2}{3}$}
		{$y=-2x-10$,  $y=-2x+\dfrac{2}{3}$}
	 \loigiai{
		Ta có  $f'(x)=x^2-4x+1$.\\
		Giả sử $M_0(x_0;y_0)$  là tiếp điểm.\\
		Hệ số góc của tiếp tuyến  tại $M_0(x_0;y_0)$ là $f'(x_0)=x_0^2-4x_0+1$.\\
		Hệ số  góc của đường thẳng $d\colon y=-2x+\dfrac{10}{3}$ là $-2$.\\
		 Tiếp tuyến song song với đường thẳng $d$ nên
		\begin{align*}
			x_0^2-4x_0+1=-2\Leftrightarrow  x_0^2-4x_0+3=0\Leftrightarrow x_0^2-4x_0+3=0\Leftrightarrow\hoac{&x=1\\ &x=3.}
		\end{align*}
		\begin{enumerate}[{\bfseries TH1:}]
			\item $x_0=1$, $y_0=\dfrac{4}{3}$, $f'(x_0)=-2$.  Phương trình tiếp tuyến là $y=-2x+\dfrac{10}{3}$.\\
			Trường hợp này loại, do trùng với đường  thẳng $d$.
			\item $x_0=3$, $y_0=-4$, $f'(x_0)=-2$.  Phương trình tiếp tuyến: $y=-2x+2$.
		\end{enumerate}
		Vậy phương trình tiếp tuyến cần tìm là  $y=-2x+2$.
	} 
\end{ex}

% nguồn: Ok_Toan11_TracNghiem_Lan2
\begin{ex}%[1D7H2-4]%[ma:fc88d413]
	Cho hàm số $y=-x^3+3x^2-3$ có đồ thị $(C)$.  Số tiếp tuyến của $(C)$ vuông góc với đường thẳng $y=-\dfrac{1}{9}x+2017$ là
	\choice 
		{\True  $2$}
		{$1$} 
		{$0$} 
		{$3$} 
	\loigiai{ 
		Gọi  $(x_0;y_0)$ là tọa độ tiếp điểm.\\
		Ta có  $y'=-3x^2+6x$.\\
		Vì tiếp tuyến của $(C)$ vuông góc  với đường thẳng $y=\dfrac{1}{9}x+2017$ nên
		\begin{align*}
			\dfrac{1}{9}\cdot y'(x_0)=-1\Leftrightarrow y'(x_0)=-9\Leftrightarrow  -3x_0^2+6x_0+9=0\Leftrightarrow\hoac{& x_0=-1\\ &x_0=3.}
		\end{align*}
		\begin{itemize}
			\item Với $x_0=-1\Rightarrow y_0=1$, suy ra phương  trình tiếp tuyến là $y=-9(x+1)+1\Leftrightarrow y=-9x-8$.
			\item Với $x_0=3\Rightarrow y_0=-3$, suy ra phương trình tiếp tuyến là  $y=-9(x-3)-3\Leftrightarrow y=-9x+24$.
		\end{itemize}
		 Vậy có $2$ tiếp tuyến thỏa mãn yêu cầu bài toán.
	} 
\end{ex}

% nguồn: Ok_Toan11_TracNghiem_Lan2
\begin{ex}%[1D7H2-4]%[ma:bf679c99]
	Cho hàm số $y=x^3-3x^2+2$ có đồ thị $(C)$. Phương trình tiếp tuyến của đồ thị $(C)$ biết tiếp tuyến song song với đường thẳng $d\colon 9x-y+7=0$  là
	 \choice
		{$y=9x+25$} 
		 {$y=-9x-25$}
		 {\True $y=9x-25$}
		 {$y=-9x+25$}
	\loigiai{ 
		Gọi  $\Delta$ là tiếp tuyến của đồ thị $(C)$ và $(x_0;y_0)$ là tọa độ tiếp điểm.\\
		 $y'=3x^2-6x$.\\
		Theo giả thiết,  $\Delta$ song song với $d\colon 9x-y+7=0\Leftrightarrow y=9x+7$ nên
		\begin{align*}
			k_\Delta=k_d=9=y'(x_0)\Leftrightarrow 3x_0^2-6x_0=9\Leftrightarrow\hoac{&  x_0=-1\\ &x_0=3.}
		\end{align*}
		\begin{itemize}
			\item Với $x_0=-1\Rightarrow y_0=-2$, ta  được $\Delta\colon y=9(x+1)-2=9x+7$ (loại).
			\item Với $x_0=3\Rightarrow y_0=2$,  ta được $\Delta\colon y=9(x-3)+2=9x-25$.
		\end{itemize}
		 Vậy tiếp tuyến cần tìm là $y=9x-25$.
	 }
\end{ex}

% nguồn: Ok_Toan11_TracNghiem_Lan2
\begin{ex}%[1D7H2-4]%[ma:9bdc2080]
	Cho đường cong $(C) \colon\, y=x^{4}-3x^{3}+2x^{2}-1$. Có bao nhiêu tiếp tuyến  của đường cong $(C)$ có hệ số góc bằng $7$?
	 \choice
		{$3$} 
		 {$2$}
		{\True $1$} 
		{$4$} 
	 \loigiai{
		\begin{itemize}
			\item Ta có $y'=4x^{3}-9x^{2}+4x$. 
			\item Hoành độ tiếp điểm là nghiệm của phương  trình $4x^{3}-9x^{2}+4x=7$.
			\item Phương trình có $1$ nghiệm nên có $1$ tiếp tuyến có hệ  số góc bằng $7$.
		\end{itemize}
	} 
\end{ex}

% nguồn: Ok_Toan11_TracNghiem_Lan2
\begin{ex}%[1D7H2-4]%[ma:9a9472eb]
	Viết phương  trình tiếp tuyến của đồ thị hàm số $y=x^3-3x^2+1$ biết nó song song với đường thẳng $y=9x+6$.
	 \choice
		{$y=9x+6$, $y=9x-6$} 
		 {\True $y=9x-26$}
		 {$y=9x+26$}
		{$y=9x-26$, $y=9x+6$} 
	\loigiai{ 
		Ta có  $y'=3x^2-6x$.\\
		 Gọi hoành độ tiếp điểm của tiếp tuyến $\Delta$ là $x_0$.\\
		Tiếp  tuyến $\Delta$ của đồ thị hàm số $y=x^3-3x^2+1$ song song với đường thẳng $y=9x+6$ nên
		\begin{align*}
			y'(x_0)=9\Leftrightarrow  3x_0^2-6x_0=9\Leftrightarrow 3x_0^2-6x_0-9=0\Leftrightarrow x_0^2-2x_0-3=0\Leftrightarrow\hoac{& x_0=-1\\ &x_0=3.}
		\end{align*}
		\begin{itemize}
			\item Với $x_0=-1\Rightarrow y(-1)=-3$, phương trình tiếp tuyến là $y=9(x+1)-3=9x+6$ (loại). 
			\item Với $x_0=3\Rightarrow y(3)=1$, phương trình tiếp  tuyến là $y=9(x-3)+1=9x-26$.
		\end{itemize}
		Vậy  có $1$ tiếp tuyến thỏa mãn yêu cầu bài toán là $y=9x-26$.
	} 
\end{ex}

% nguồn: Ok_Toan11_TracNghiem_Lan2
\begin{ex}%[1D7H2-4]%[ma:dfbea4e7]
	Có bao nhiêu tiếp tuyến của đồ thị hàm  số $y=-x^3+2x^2$ song song với đường thẳng $y=x$?
	\choice 
		{$3$} 
		{$2$} 
		{$0$} 
		{\True  $1$}
	\loigiai{ 
		Giả sử tiếp tuyến của đồ thị hàm  số $y=-x^3+2x^2$ tại $M(x_0,y_0)$ có dạng: $y=y'(x_0)(x-x_0)+y_0$.\\
		 Vì tiếp tuyến song song với đường thẳng $y=x$ nên
		\begin{align*}
			y'(x_0)=1\Leftrightarrow-3x_0^2+4x_0=1\Leftrightarrow 3x_0^2-4x_0+1=0\Leftrightarrow\hoac{&  x_0=1\\ &x_0=\dfrac{1}{3}.}
		\end{align*}
		\begin{itemize}
			 \item Với $x_0=1$, $y_0=1$, phương trình tiếp tuyến là $y=x$ (loại).
			\item Với $x_0=\dfrac{1}{3}$, $y_0=\dfrac{5}{27}$, phương trình  tiếp tuyến là $y=x-\dfrac{4}{27}$.
		\end{itemize}
		Vậy  có một tiếp tuyến thỏa mãn yêu cầu bài toán.
	 }
\end{ex}

% nguồn: Ok_Toan11_TracNghiem_Lan2
\begin{ex}%[1D7H2-4]%[ma:dca2e47d]
	Cho hàm số $y=x^{3}-3x^{2}+2$ có  đồ thị. Tìm số tiếp tuyến của đồ thị song song với đường thẳng $d \colon\, y=9x-25$.
	\choice 
		{\True  $1$}
		 {$2$}
		{$3$} 
		 {$0$}
	 \loigiai{
		\begin{itemize}
			\item Ta có hàm số $y=x^{3}-3x^{2}+2$,  có $y'=3x^{2}-6x$.
			\item Gọi $M\left(x_{0};y_{0}\right)$ là tọa độ tiếp điểm của tiếp tuyến và  đồ thị $(C)$, khi đó hệ số góc của tiếp tuyến là $k=3x_{0}{}^{2}-6x_{0}$.
			\item Tiếp tuyến của  $(C)$ song song với đường thẳng $y=9x-25$ nên ta có
			\begin{eqnarray*}
				k=9\Leftrightarrow3{x_{0}}^{2}-6x_{0}=9 \Leftrightarrow \hoac{& x_{0}=-1\, \left(y_{0}=-2\right) \\ & x_{0}=3\, \left(y_{0}=2\right).} 
			\end{eqnarray*}
			\item Với $M\left(-1;-2\right)$ phương trình tiếp tuyến của $(C)$ là $y=9x+7$. 
			\item Với  $M\left(3;2\right)$ phương trình tiếp tuyến của $(C)$ là $y=9x-25$.
			\item Tiếp tuyến $y=9x-25$ trùng với đường thẳng  đã cho nên bị loại.
			\item Vậy có $1$ tiếp tuyến cần tìm là $y=9x+7$. 
		\end{itemize}
	 }
\end{ex}

% nguồn: Ok_Toan11_TracNghiem_Lan2
\begin{ex}%[1D7H2-4]%[ma:08ab0c93]
	Phương trình tiếp tuyến của đồ thị hàm số $(C)\colon y=\dfrac{2x+1}{x+2}$ song song với đường thẳng $d\colon  y=3x+2$ là
	 \choice
		 {$y=3x+2$}
		{$y=3x-2$} 
		{\True $y=3x+14$} 
		{$y=3x+5$} 
	 \loigiai{
		Vì tiếp tuyến của đồ thị $(C)$ song song với $\Delta\colon y=3x+2$ nên gọi toạ độ tiếp điểm  là $M(x_0;y_0)$ ta có
		\begin{align*}
			y'(x_0)=3\Leftrightarrow\dfrac{3}{(x_0+2)^2}=3\Leftrightarrow(x_0+2)^2=1\Leftrightarrow\hoac{& x_0=-1\\ &x_0=-3.} 
		\end{align*}
		\begin{itemize}
			 \item Với $x_0=-1\Rightarrow y_0=-1$, ta được $\Delta\colon y=3(x+1)-1=3x+2$ (loại).
			\item Với  $x_0=-3\Rightarrow y_0=5$, ta được $\Delta\colon y=3(x+3)+5=3x+14$.
		\end{itemize}
		Vậy  có một tiếp tuyến song song với đường thẳng $d\colon y=3x+2$ có phương trình $y=3x+14$.
	} 
\end{ex}

% nguồn: Ok_Toan11_TracNghiem_Lan2
\begin{ex}%[1D7H2-4]%[ma:af2813bf]
	Cho  hàm số $y=x^3-3x$ có đồ thị $(C)$. Số các tiếp tuyến của đồ thị $(C)$ song song với đường thẳng $y=3x-10$ là
	\choice 
		{\True $2$} 
		 {$1$}
		{$3$} 
		 {$0$}
	 \loigiai{
		Ta có  $y=x^3-3x\Rightarrow y'=3x^2-3$.\\
		Gọi $M(x_0;y_0)$ là  tiếp điểm.\\
		Vì tiếp tuyến  song song với đường thẳng $y=3x-10$ nên
		\begin{align*}
			f'(x_0)=3\Leftrightarrow 3x_0^2-3=3\Leftrightarrow  x_0^2=2\Leftrightarrow x_0=\pm\sqrt{2}.
		\end{align*}
		\begin{itemize}
			 \item Với $x_0=\sqrt{2}\Rightarrow y_0=-\sqrt{2}$. Khi đó, phương trình tiếp tuyến là
			\begin{align*}
				y=3\left(x-\sqrt{2}\right)-\sqrt{2}=3x-4\sqrt{2}. 
			\end{align*}
			 \item Với $x_0=-\sqrt{2}\Rightarrow y_0=\sqrt{2}$. Khi đó, phương trình tiếp tuyến là
			\begin{align*}
				y=3\left(x+\sqrt{2}\right)+\sqrt{2}=3x+4\sqrt{2}. 
			\end{align*}
		\end{itemize}
		Vậy có $2$  tiếp tuyến thỏa mãn yêu cầu bài toán.
	} 
\end{ex}

% nguồn: Ok_Toan11_TracNghiem_Lan2
\begin{ex}%[1D7H2-4]%[ma:2942ab72]
	Có bao nhiêu điểm $M$ thuộc đồ thị hàm số $f(x)=x^3+1$ sao cho tiếp tuyến của đồ thị  hàm số $f(x)$ tại $M$ song song với đường thẳng $d\colon y=3x-1$?
	 \choice
		 {$3$}
		 {$2$}
		 {$0$}
		 {\True $1$}
	 \loigiai{
		Gọi $M(a;a^3+1)$ là điểm  thuộc đồ thị $(C)\colon f(x)=x^3+1$.\\
		Ta có $f '(x)=3x^2$. Suy ra phương  trình tiếp tuyến của $(C)$ tại $M$ là
		\begin{align*}
			 y=3a^2(x-a)+a^3+1\Leftrightarrow y=3a^2x-2a^3+1.
		\end{align*}
		 Gọi $\Delta\colon y=3a^2x-2a^3+1$. Khi đó,
		\begin{align*}
			\Delta \parallel  d\Leftrightarrow\heva{& 3a^2=3\\ &-2a^3+1\neq-1}\Leftrightarrow\heva{& a=\pm 1\\ &a\neq 1}\Leftrightarrow a=-1.
		\end{align*}
		Vậy có duy  nhất điểm $M$ thỏa mãn yêu cầu là $M(-1;0)$.
	 }
\end{ex}

% nguồn: Ok_Toan11_TracNghiem_Lan2
\begin{ex}%[1D7H2-4]%[ma:5de2a488]
	Cho hàm số $y=\dfrac{x^3}{3}+3x^2-2$ có  đồ thị là $(C)$. Viết phương trình tiếp tuyến với đồ thị $(C)$ biết tiếp tuyến có hệ số góc $k=-9$.
	\choice 
		 {$y+16=-9(x+3)$}
		{$y=-9(x+3)$} 
		 {$y-16=-9(x-3)$}
		 {\True $y-16=-9(x+3)$}
	 \loigiai{
		Ta  có $y'=x^2+6x$.\\
		Gọi $M(x_0;y_0)$ là  tiếp điểm.\\
		Do tiếp tuyến có hệ số góc  $k=-9$ nên
		\begin{align*}
			 y'(x_0)=-9\Leftrightarrow x_0^2+6x_0+9=0\Leftrightarrow x_0=-3.
		\end{align*}
		 Khi đó, $y_0=y(-3)=16$.\\
		Vậy $y=y'(x_0)(x-x_0)+y_0=-9(x+3)+16$ hay  $y-16=-9(x+3)$.
	 }
\end{ex}

% nguồn: Ok_Toan11_TracNghiem_Lan2
\begin{ex}%[1D7H2-4]%[ma:d105e297]
	Tìm các tiếp tuyến của đồ thị hàm số $y=\dfrac{2x+1}{x-1}$ biết các tiếp tuyến đó song song  với đường thẳng $y=-3x$.
	 \choice
		 {\True $y=-3x+11$, $y=-3x-1$}
		{$y=-3x-6$,  $y=-3x-11$}
		{$y=-3x+1$} 
		 {$y=-3x+6$}
	 \loigiai{
		\begin{itemize}
			\item  Ta có $y'=\dfrac{-3}{(x-1)^{2}}$.
			\item Gọi $\Delta$ là  tiếp tuyến cần tìm.
			\item Tiếp tuyến $\Delta$ song  song với đường thẳng $y=-3x$ suy ra hệ số góc của tiếp tuyến $\Delta$ là $k=-3$.
			\item Gọi $M_{0}\left(x_{0};y_{0}\right)$ là tiếp điểm, ta có $y'\left(x_{0}\right)=k  \Leftrightarrow \dfrac{-3}{(x_{0}-1)^{2}}=-3 \Leftrightarrow \hoac{& x_{0}=2\, \left(y_{0}=5\right) \\ & x_{0}=0\, \left(y_{0}=-1\right).}$.
			\item Phương trình tiếp tuyến của đồ thị hàm số đã cho tại  điểm $M_{1}(2;5)$ là $y=-3(x-2)+5 \Leftrightarrow y=-3x+11$.
			\item Phương  trình tiếp tuyến của đồ thị hàm số đã cho tại điểm $M_{2}(0;-1)$ là $y=-3(x-0)-1 \Leftrightarrow y=-3x-1$.
			\item Vậy có $2$  tiếp tuyến cần tìm là $y=-3x+11$ và $y=-3x-1$.
		\end{itemize}
	 }
\end{ex}

% nguồn: Ok_Toan11_TracNghiem_Lan2
\begin{ex}%[1D7H2-4]%[ma:f53d3070]
	Cho hàm số $y=x^3-3x^2+1$ có  đồ thị là $(C)$. Phương trình các tiếp tuyến của $(C)$ song song với đường thẳng $y=9x+10$ là
	 \choice
		{$y=9x+6$, $y=9x-28$} 
		{$y=9x$, $y=9x-26$} 
		{$y=9x-6$,  $y=9x-28$}
		 {\True $y=9x+6$, $y=9x-26$}
	\loigiai{ 
		Ta có  $y'=3x^2-6x$.\\
		Hệ  số góc của tiếp tuyến là
		\begin{align*}
			k=y'(x_0)=3x_0^2-6x_0=9\Leftrightarrow 3x_0^2-6x_0-9=0\Leftrightarrow x_0^2-2x_0-3=0\Leftrightarrow\hoac{& x_0=-1\\  &x_0=3.}
		\end{align*}
		Phương trình tiếp tuyến tại  $M(3;1)$ là $y=9(x-3)+1=9x-26$.\\
		Phương trình tiếp tuyến tại $N(-1;-3)$ là  $y=9(x+1)-3=9x+6$.
	} 
\end{ex}

% nguồn: Ok_Toan11_TracNghiem_Lan2
\begin{ex}%[1D7H2-4]%[ma:83d6a8dd]
	Số tiếp tuyến của đồ thị hàm số $y=-x^4+2x^2$ song song với trục  hoành là
	\choice 
		{$3$} 
		 {$2$}
		{$0$} 
		 {\True $1$}
	\loigiai{ 
		Ta có $y'=-4x^3+4x$.\\ 
		 Gọi $M(x_0;y_0)$ là tiếp điểm.\\
		Vì tiếp tuyến song song với  trục hoành nên có hệ số góc bằng $0$.\\
		Suy ra  $y'(x_0)=0\Leftrightarrow-4x_0^3+4x_0=0\Leftrightarrow\hoac{& x_0=0\\ &x_0=-1\\ &x_0=1.}$
		\begin{itemize}
			\item Với $x_0=0$ thì $y_0=0$, tiếp tuyến là $y=0$ (loại). 
			\item Với  $x_0=-1$ thì $y_0=1$, tiếp tuyến là $y=1$.
			 \item Với $x_0=1$ thì $y_0=1$, tiếp tuyến là $y=1$.
		\end{itemize}
		Vậy có một tiếp tuyến song song với  trục hoành có phương trình $y=1$.
	} 
\end{ex}

% nguồn: Ok_Toan11_TracNghiem_Lan2
\begin{ex}%[1D7H2-4]%[ma:0d8cf18a]
	 Cho hàm số $y=\dfrac{2x-1}{x+1}$ có đồ thị $(C)$. Tiếp tuyến của $(C)$ vuông góc với đường thẳng $x+3y+2=0$ tại điểm có hoành độ
	\choice 
		{$x=0$} 
		{$x=-2$} 
		{\True  $\hoac{& x=0\\ & x=-2}$}
		{$\hoac{& x=0\\ & x=2}$} 
	\loigiai{ 
		Tiếp tuyến của $(C)$ vuông góc với đường thẳng $x+3y+2=0\Leftrightarrow  y=-\dfrac{1}{3}x-\dfrac{2}{3}$ nên hệ số góc của tiếp tuyến là $k=3$.\\
		Hoành độ tiếp điểm là nghiệm của  phương trình
		\begin{align*}
			y'=3\Leftrightarrow\dfrac{3}{(x+1)^2}=3\Leftrightarrow(x+1)^2=1\Leftrightarrow\hoac{& x=0\\ & x=-2.} 
		\end{align*}
		Vậy hoành độ tiếp điểm cần tìm là  $\hoac{& x=0\\ & x=-2.}$
	} 
\end{ex}

% nguồn: Ok_Toan11_TracNghiem_Lan2
\begin{ex}%[1D7H2-4]%[ma:1f7c1228]
	Cho hàm số  $f(x)=x^3-3x^2$. Tiếp tuyến của đồ thị hàm số song song với đường thẳng $y=9x+5$ là
	\choice 
		 {$y=9(x+3)$}
		{\True $y=9(x-3)$} 
		 {$y=9x+5$ và $y=9(x-3)$}
		 {$y=9x+5$}
	\loigiai{ 
		Ta có $f'(x)=3x^2-6x$.\\ 
		Tiếp tuyến song song  với đường thẳng $y=9x+5$ nên $3x^2-6x=9\Leftrightarrow\hoac{&x=-1\\ &x=3.}$
		\begin{itemize}
			 \item Với $x=-1\Rightarrow y=-4$, $f'(-1)=9$. Phương trình tiếp tuyến là $y=9x+5$ (loại).
			\item Với $x=3\Rightarrow y=0$, $f '(3)=9$.  Phương trình tiếp tuyến là $y=9(x-3)$.
		\end{itemize}
		Vậy tiếp tuyến  cần tìm là $y=9(x-3)$.
	 }
\end{ex}

% nguồn: Ok_Toan11_TracNghiem_Lan2
\begin{ex}%[1D7H2-4]%[ma:c9ba4d20]
	Tìm điểm $M$ có hoành độ âm trên đồ thị $(C) \colon\, y=\dfrac{1}{3}x^{3}-x+\dfrac{2}{3}$ sao cho tiếp tuyến tại  $M$ vuông góc với đường thẳng $y=-\dfrac{1}{3}x+\dfrac{2}{3}$.
	\choice 
		{$M\left(-1;\dfrac{4}{3}\right)$} 
		{\True  $M\left(-2;0\right)$}
		{$M\left(2;\dfrac{4}{3}\right)$} 
		 {$M\left(-2;-4\right)$}
	 \loigiai{
		\begin{itemize}
			\item Ta có $y'(x)=x^{2}-1$. 
			\item Tiếp tuyến tại $M$ vuông góc với đường thẳng $y=-\dfrac{1}{3}x+\dfrac{2}{3}$ nên tiếp tuyến có hệ  số góc $k=3$.
			\item Xét phương trình $y'(x)=3 \Leftrightarrow x^{2}-1=3 \Leftrightarrow x^{2}=4 \Leftrightarrow \hoac{& x=2 \\ & x=-2.}$. 
			\item Do $M$ có hoành độ âm nên $x=-2$  thỏa mãn, $x=2$ loại.
			\item Với $x=-2$ thay vào phương trình $(C)$,  ta được $y=0$.
			\item Vậy điểm $M$ cần  tìm là $M\left(-2;0\right)$.
		\end{itemize}
	 }
\end{ex}

% nguồn: Ok_Toan11_TracNghiem_Lan2
\begin{ex}%[1D7H2-4]%[ma:fab01ff6]
	Có bao nhiêu tiếp tuyến  của đồ thị hàm số $y=-x^3+2x^2$ song song với đường thẳng $y=x$?
	 \choice
		{$2$} 
		 {$4$}
		{$3$} 
		{\True  $1$}
	 \loigiai{
		 Gọi $M(x_0;y_0)$ là tiếp điểm của tiếp tuyến song song với đường thẳng $y=x$ của đồ thị hàm số $y=-x^3+2x^2$.\\
		 Khi đó, ta có
		\begin{align*}
			y'(x_0)=1\Leftrightarrow-3x_0^2+4x_0=1\Leftrightarrow 3x_0^2-4x_0+1=0\Leftrightarrow\hoac{&  x_0=1\\ &x_0=\dfrac{1}{3}.}
		\end{align*}
		\begin{itemize}
			\item  Với $x_0=1$ ta được $M(1;1)$, phương trình tiếp tuyến: $y=1\cdot (x-1)+1\Leftrightarrow y=x$ (loại).
			 \item Với $x_0=\dfrac{1}{3}$ ta được $M\left(\dfrac{1}{3};\dfrac{5}{27}\right)$, phương trình tiếp tuyến: $y=1\cdot \left(x-\dfrac{1}{3}\right)+\dfrac{5}{27}=x-\dfrac{4}{27}$.
		\end{itemize}
		Vậy chỉ có một tiếp  tuyến thỏa mãn yêu cầu bài toán.
	} 
\end{ex}

% nguồn: Ok_Toan11_TracNghiem_Lan2
\begin{ex}%[1D7H2-4]%[ma:30c53ba0]
	Cho hàm số $f(x)=\dfrac{2x+1}{x-1}$ $(C)$. Các tiếp tuyến của $(C)$ song song với đường  thẳng $y=-3x$ có phương trình là
	 \choice
		{\True $y=-3x-1$;  $y=-3x+11$}
		{$y=-3x+10$; $y=-3x-4$} 
		 {$y=-3x+5$; $y=-3x-5$}
		 {$y=-3x+2$; $y=-3x-2$}
	 \loigiai{
		Gọi $M(x_0;y_0)$ là tiếp điểm của tiếp tuyến. Theo giả  thiết ta có
		\begin{align*}
			f'(x_0)=-3\Leftrightarrow\dfrac{-3}{(x_0-1)^2}=-3\Leftrightarrow(x_0-1)^2=1\Leftrightarrow\hoac{&  x_0=0\\ &x_0=2.}
		\end{align*}
		\begin{itemize}
			\item Với $x_0=0\Rightarrow y_0=-1$: Phương trình tiếp tuyến: $y=-3(x-0)-1\Leftrightarrow  y=-3x-1$.
			 \item Với $x_0=2\Rightarrow y_0=5$: Phương trình tiếp tuyến: $y=-3(x-2)+5\Leftrightarrow y=-3x+11$.
		\end{itemize}
		Ta thấy cả hai  tiếp tuyến đều thỏa mãn điều kiện đề bài.
	} 
\end{ex}

% nguồn: Ok_TracNghiemLop11_Lan1
\begin{ex}%[1D7H2-4]%[1D5B2-3]%[ma:d1dc6306]
	Cho hàm số $y=x^2+mx+1$ có đồ thị $C_m$. Xác định $m$ sao cho tiếp  tuyến với đồ thị $(C_m) $ tại điểm có hoành độ $x=1$ song song với đường thẳng $y=3x-17$.
	\choice 
		{\True  $m=1$}
		 {$m=-1$}
		{$m=3$} 
		 {$m=0$}
	 \loigiai{
		Ta có $y'=2x+m$. Theo bài, tiếp tuyến của đồ thị $(C_m)$ tại điểm có hoành độ $x=1$ song song với đường thẳng $y=3x-17$ nên suy ra $y'(1)=3 \Leftrightarrow 2+m=3 \Leftrightarrow m=1$. Với $m=1$, ta có phương trình tiếp tuyến  là $y=3(x-1)+3\Leftrightarrow y=3x$ (thỏa mãn).
	 }
\end{ex}

% nguồn: Ok_TracNghiemLop11_Lan1
\begin{ex}%[1D7H2-4]%[ma:c868d4a9]
	Cho hàm số $y=\dfrac{x-3}{x+2}$. Biết tiếp tuyến đồ thị hàm số vuông góc với đường thẳng $y=-\dfrac{1}{5}x+2017$. Tìm tất  cả các hoành độ tiếp điểm $x_0$.
	 \choice
		{$x_0=-1$} 
		{$x_0=-1$; $x_0=2$} 
		{\True $x_0=-1$;  $x_0=-3$}
		{$x_0=3$} 
	\loigiai{ 
	Tập xác  định của hàm số $\mathscr{D} = \mathbb{R}\setminus\{-2\}$.\\
	Ta có $y' = \dfrac{5}{(x+2)^2}$.\\ 
	 Gọi $x_0$ là hoành độ tiếp điểm.\\
	Vì tiếp tuyến của đồ thị hàm số vuông góc  với đường thẳng $y=-\dfrac{1}{5}x+2017$ nên
	\begin{align*}
		\dfrac{5}{(x_0+2)^2} = 5 \Leftrightarrow\heva{& x_0+2 \neq 0 \\& (x_0+2)^2 = 1} \Leftrightarrow \heva{& x_0 \neq -2 \\& \hoac{& x_0=-1 \\& x_0=-3}} \Leftrightarrow \hoac{& x_0 = -1 \\&  x_0 = -3.}
	\end{align*}
	 }
\end{ex}

% nguồn: Ok_Toan11_TracNghiem_Lan2
\begin{ex}%[1D7V2-4]%[ma:f1b53bc1]
	Đường thẳng $y=6x+m+1$ là tiếp tuyến của đồ  thị hàm số $y=x^{3}+3x-1$ khi $m$ bằng
	 \choice
		{$-4$  hoặc $-2$}
		{\True $-4$ hoặc  $0$}
		{$0$  hoặc $2$}
		{$-2$ hoặc $2$} 
	 \loigiai{
		\begin{itemize}
			\item Gọi $(C)$ là đồ thị  hàm số $y=x^{3}+3x-1$.
			\item Ta có  $y'=3x^{2}+3$.
			\item Hệ số  góc của tiếp tuyến là $k=6$. Ta có
			\begin{eqnarray*}
				y'=6 \Leftrightarrow  3x^{2}+3=6 \Leftrightarrow \hoac{& x=1 \Rightarrow y=3 \\ & x=-1 \Rightarrow y=-5.}
			\end{eqnarray*}
			\item Phương trình  tiếp tuyến của $(C)$ tại điểm $M\left(1;3\right)$ là $y=6x-3$.
			\item Phương trình  tiếp tuyến của $(C)$ tại điểm $N\left(-1;-5\right)$ là $y=6x+1$.
			\item Để đường thẳng $y=6x+m+1$ là tiếp tuyến của $(C)$ thì $\hoac{& m+1=-3 \\ &  m+1=1.} \Leftrightarrow \hoac{& m=-4 \\ & m=0.}$
		\end{itemize}
	 }
\end{ex}

% nguồn: Ok_Toan11_TracNghiem_Lan2
\begin{ex}%[1D7V2-4]%[ma:c611352e]
	Cho hàm  số $y=x^{3}-3x^{2}+2$ có đồ thị $(C)$. Tìm số tiếp tuyến của đồ thị $(C)$ song song với đường thẳng $d \colon\, y=9x-25$.
	\choice 
		{\True $1$} 
		 {$3$}
		 {$0$}
		 {$2$}
	 \loigiai{
		\begin{itemize}
			\item Ta có  $y'=3x^{2}-6x$.
			\item Vì  tiếp tuyến của $(C)$ song song với đường thẳng $d \colon\, y=9x-25$ nên có
			\begin{eqnarray*}
				y'=9\Leftrightarrow 3x^{2}-6x=9 \Leftrightarrow x^{2}-2x-3=0 \Leftrightarrow \hoac{& x=-1\, \left(y=-2\right) \\  & x=3\, \left(y=2\right) .}
			\end{eqnarray*}
			\item Phương trình tiếp tuyến của đồ thị hàm số đã cho tại điểm $M(-1;-2)$ là  $y=9(x+1)-2 \Leftrightarrow y=9x+7$.
			 \item Phương trình tiếp tuyến của đồ thị hàm số đã cho tại điểm $M(3;2)$ là $y=9(x-3)+2 \Leftrightarrow y=9x-25$.
			\item Tiếp tuyến $y=9x-25$ trùng với $d$  nên bị loại.
			\item Vậy chỉ có $1$  tiếp tuyến thỏa mãn yêu cầu bài toán.
		\end{itemize}
	 }
\end{ex}

% nguồn: Ok_Toan11_TracNghiem_Lan2
\begin{ex}%[1D7V2-4]%[ma:aa0c2ea5]
	Cho hàm số $y = x^4-2x^2+m-2$ có  đồ thị $(C)$. Gọi $S$ là tập các giá trị của $m$ sao cho đồ thị $(C)$ có đúng một tiếp tuyến song song với trục $Ox$. Tổng các phần tử của $S$ là
	\choice 
		 {$3$}
		{$8$} 
		{\True  $5$}
		 {$2$}
	 \loigiai{
		\begin{itemize}
			\item Vì tiếp tuyến song song với trục $Ox$ nên hệ số góc của tiếp tuyến  $k=0$.
			\item Gọi tiếp điểm là $M\left(x_{0};y_{0}\right) \in (C)$,  khi đó $y'\left(x_{0}\right)=4{x_{0}}^{3}-4x_{0}=0 \Leftrightarrow \hoac{& x_{0}=0\, \left(y_{0}=m-2\right) \\ & x_{0}=\pm1\, \left(y_{0}=m-3\right).}$.
			\item Để có đúng một tiếp tuyến song song với trục $Ox$, cần có một giá trị tung độ $y_0 \neq 0$ và giá trị còn lại  bằng $0$ (trùng $Ox$).
			\item TH1: $m-2 \neq 0$  và $m-3=0 \Leftrightarrow m=3$.
			\item TH2:  $m-2 = 0$ và $m-3 \neq 0 \Leftrightarrow m=2$.
			\item  Vậy $S=\{3;2\}$.
			\item Suy ra tổng các giá trị các phần tử của  $S$ là $3+2=5$.
		\end{itemize}
	} 
\end{ex}

% nguồn: Ok_Toan11_TracNghiem_Lan2
\begin{ex}%[1D7V2-4]%[ma:b8d600c8]
	Tiếp tuyến của đồ thị hàm số $y=2x^{3}-3x^{2}-12x+1$ song song với đường thẳng $d \colon\, 12x+y=0$ có dạng là $y=ax+b$. Tính giá trị của $2a+b$. 
	\choice 
		{$ -23 $ hoặc  $ -24 $}
		{\True  $-23$}
		{$-24$} 
		{$0$} 
	\loigiai{ 
		\begin{itemize}
			\item Ta có $d \colon\,  12x+y=0 \Rightarrow y=-12x$. Hệ số góc của đường thẳng $d$ là $k_{d}=-12$.
			\item Do tiếp tuyến của đồ thị hàm số song song với đường thẳng $d$ nên hệ số góc  của tiếp tuyến là $k_{tt}=k_{d}=-12$.
			\item Ta có  $y'=6x^{2}-6x-12$.
			\item Giả sử $M\left(x_{0};y_{0}\right)$ là tiếp điểm. Khi đó $k_{tt}=y'\left(x_{0}\right)=6{x_{0}}^{2}-6x_{0}-12$.  Suy ra
			\begin{eqnarray*}
				k_{tt}=-12 \Leftrightarrow  6{x_{0}}^{2}-6x_{0}-12=-12\Leftrightarrow \hoac{& x_{0}=0\, \left(y_{0}=1\right) \\ & x_{0}=1\, \left(y_{0}=-12\right).}
			\end{eqnarray*}
			\item Tiếp  tuyến tại $M(0;1)$ là $y=-12(x-0)+1\Leftrightarrow y=-12x+1$.
			\item Tiếp tuyến  tại $N(1;-12)$ là $y=-12(x-1)-12\Leftrightarrow y=-12x$. (Loại vì trùng $d$)
			\item Vậy $y=-12x+1$, suy ra  $a=-12$, $b=1$. Suy ra $2a+b=-23$.
		\end{itemize}
	 }
\end{ex}

% nguồn: Ok_Toan11_TracNghiem_Lan2
\begin{ex}%[1D7C2-4]%[ma:11ee9a5b]
	Cho hàm số $y=\dfrac{x+2}{2x+3}$. Đường thẳng $d \colon\, y=ax+b$ là tiếp tuyến của đồ thị hàm số. Biết $d$ cắt trục hoành, trục tung lần lượt tại hai điểm $A$, $B$ sao  cho $\Delta OAB$ cân tại $O$. Khi đó $a+b$ bằng
	 \choice
		{$-1$} 
		 {$0$}
		 {$2$}
		 {\True $-3$}
	 \loigiai{
		\begin{itemize}
			\item  Tập xác định của hàm số $y=\dfrac{x+2}{2x+3}$ là $\mathscr{D}=\mathbb{R}\setminus\left\{-\dfrac{3}{2}\right\}$.
			\item Hệ số góc tiếp tuyến là  $k = \pm \dfrac{OB}{OA} = \pm-1$.
			\item Gọi $\left(x_{0};y_{0}\right)$ là tọa độ tiếp điểm. 
			\item Mặt khác,  $k=y'\left(x_{0}\right)=-\dfrac{1}{(2x_{0}+3)^{2}}<0$. Suy ra $-\dfrac{1}{(2x_{0}+3)^{2}}=-1\Leftrightarrow \hoac{& x_0=-2\, \left(y_{0}=0\right)\\ & x_0=-1\, \left(y_{0}=1\right).}$
			\item Phương trình tiếp tuyến  tại điểm $M(-2;0)$ là $y=-x-2$. (Thỏa mãn).
			\item  Phương trình tiếp tuyến tại điểm $N(-1;1)$ là $y=-x$. (Loại vì $A$, $B$ trùng với $O$)
			\item Vậy $a=-1$, $b=-2$. Suy  ra $a+b=-3$.
		\end{itemize}
	} 
\end{ex}

% nguồn: Ok_Toan11_TracNghiem_Lan2
\begin{ex}%[1D7C2-4]%[ma:cda3d585]
	Cho hàm số $y=\dfrac{2x-1}{x-1}$  có đồ thị. Có bao nhiêu tiếp tuyến của đồ thị cắt trục $Ox$, $Oy$ lần lượt tại hai điểm $A$ và $B$ thỏa mãn điều kiện $OA=4OB$.
	\choice 
		{\True $2$} 
		{$3$} 
		{$1$} 
		 {$4$}
	\loigiai{ 
		\begin{itemize}
			 \item Tập xác định của hàm số $y=\dfrac{2x-1}{x-1}$ là $\mathscr{D}=\mathbb{R}\setminus\{1\}$.
			\item  Hệ số góc tiếp tuyến là $k = \pm \dfrac{OB}{OA} = \pm \dfrac{1}{4}$.
			\item Gọi  $\left(x_{0};y_{0}\right)$ là tọa độ tiếp điểm.
			\item Mặt khác, $k=y'\left(x_{0}\right)=-\dfrac{1}{(x_{0}-1)^{2}}<0$. Suy ra $-\dfrac{1}{(x_{0}-1)^{2}}=-\dfrac{1}{4}\Leftrightarrow \hoac{&  x_0=3\, \left(y_{0}=\dfrac{5}{2}\right) \\ & x_0=-1\, \left(y_{0}=\dfrac{3}{2}\right).}$
			 \item Phương trình tiếp tuyến tại điểm $M\left(3;\dfrac{5}{2}\right)$ là $y=-\dfrac{1}{4}x+\dfrac{13}{4}$.
			\item Phương trình tiếp  tuyến tại điểm $N\left(-1;\dfrac{3}{2}\right)$ là $y=-\dfrac{1}{4}x+\dfrac{5}{4}$.
			\item Vậy có $2$ tiếp tuyến thỏa mãn yêu cầu bài  toán.
		\end{itemize}
	} 
\end{ex}

% nguồn: Ok_Toan11_TracNghiem_Lan2
\begin{ex}%[1D7C2-4]%[ma:6d0e12a7]
	Cho hàm số $y=\dfrac{x+2}{2x+3}$ có đồ thị $(C)$. Đường thẳng $d$ có phương trình $y=ax+b$ là tiếp tuyến của $(C)$ biết $d$ cắt trục hoành tại $A$ và cắt trục tung tại $B$ sao cho tam giác  $\Delta OAB$ cân tại $O$, với $O$ là gốc tọa độ. Tính $a+b$.
	\choice 
		 {$-1$}
		 {$-2$}
		 {$0$}
		{\True  $-3$}
	\loigiai{ 
		\begin{itemize}
			\item Tập  xác định của hàm số $y=\dfrac{x+2}{2x+3}$ là $\mathscr{D}=\mathbb{R}\setminus\left\{-\dfrac{3}{2}\right\}$.
			\item  Ta có $y'=\dfrac{-1}{(2x+3)^{2}}<0$.
			\item Tam  giác $\Delta OAB$ cân tại $O$, suy ra hệ số góc của tiếp tuyến bằng $\pm1$.
			\item  Do $y'<0$, hệ số góc của tiếp tuyến là $k=-1$.
			\item Gọi tọa  độ tiếp điểm là $\left(x_{0};y_{0}\right)$ $\left(x_{0}\ne-\dfrac{3}{2}\right)$. Ta có
			\begin{eqnarray*}
				\dfrac{-1}{\left(2x_{0}+3\right)^{2}}=-1 \Rightarrow \left(2x_0+3\right)^2=1 \Leftrightarrow \hoac{&  x_{0}=-2\,\left( y_0=0\right) \\ & x_{0}=-1\, \left(y_0=1\right).}
			\end{eqnarray*}
			\item Với $x_{0}=-1$, $y_{0}=1$, phương trình tiếp tuyến là $y=-x$. (Loại vì $A$, $B$ trùng  $O$).
			\item Với $x_{0}=-2$, $y_{0}=0$, phương trình tiếp tuyến là $y=-(x+2)+0  \Leftrightarrow y=-x-2$. (Thỏa mãn).
			\item Vậy  $a=-1$, $b=-2 \Rightarrow a+b=-3$.
		\end{itemize}
	} 
\end{ex}

% nguồn: Ok_TracNghiemLop11_Lan1
\begin{ex}%[1D7K2-4]%[ma:b37fe2fd]
	Cho hàm số $y  = 4x + 2\cos 2x$ có đồ thị là $(C)$. Hoành độ của các điểm trên $(C)$ mà tại đó tiếp tuyến của $(C)$ song song hoặc trùng với trục hoành là
	 \choice
		{\True $x =  \dfrac{\pi}{4} + k\pi \enspace \left( k \in \mathbb{Z} \right)$}
		{$x = \pi + k\pi \enspace \left( k \in \mathbb{Z}  \right)$}
		{$x =  \dfrac{\pi}{2} + k\pi \enspace \left( k \in \mathbb{Z} \right)$}
		 {$x = k2\pi \enspace \left( k \in \mathbb{Z} \right)$}
	\loigiai{ 
		$y'=4 -  4\sin 2x$. \\
		Gọi $M(x_0; y_0)$ là tọa độ tiếp điểm của $(C)$ và tiếp  tuyến cần tìm.\\
		Trục hoành có phương trình $y = 0$. Để tiếp tuyến của $(C)$ song song hoặc trùng với trục hoành thì hệ số góc của tiếp tuyến cần tìm là  $y'(x_0) = 0$. Suy ra \\
		$4 - 4\sin 2x_0 = 0 \Leftrightarrow \sin 2x_0 = 1 \Leftrightarrow 2x_0 = \dfrac{\pi}{2} + k2\pi \Leftrightarrow x_0 = \dfrac{\pi}{4} +  k\pi \ (k\in\mathbb{Z})$.
	} 
\end{ex}


%%%%%%%%%%%%%%%% ĐÚNG/SAI %%%%%%%%%%%%%%%%

% nguồn: Ok_dungsai_lop11_Lan1
\begin{ex}%[1D7H2-4]%[ma:3801f31d]
 Cho hàm số $y=\dfrac{x^3}{3}+3x^2-2$ có đồ thị là  $(C)$. Khi đó, các khẳng định sau đây đúng hay sai?
 \choiceTF 
  {\True Phương trình tiếp tuyến với đồ thị $(C)$ có hệ  số góc $k=-9$ là $y-16=-9(x+3)$}
  {\True Số phương trình tiếp tuyến  có hệ số góc $k=0$ là $2$}
  {Với $k<-9$ thì đồ thị $(C)$ luôn có phương trình tiếp tuyến} 
  {Với $k>-9$ thì đồ thị $(C)$ có $1$ phương trình tiếp  tuyến}
  \loigiai{
  Gọi $M\left(x_0;\dfrac{x_0^3}{3}+3x_0^2-2\right)$ là tiếp  điểm.\\
  Ta có $k=f'(x_0) \Leftrightarrow  x_0^2+6x_0=k$ (*).
  \begin{itemize}
    \item $k=-9\Rightarrow (*) \Leftrightarrow x_0=-3\Rightarrow y_0=f(x_0)=16$.\\
   Vậy phương trình  tiếp tuyến là $y-16=-9(x+3)$.
   \item Với $k=0\Rightarrow (*) \Leftrightarrow \heva{&x=0\\&x=-6}\Rightarrow $ có $2$ phương  trình tiếp tuyến có hệ số góc $k=0$.
   \item Để đồ thị $(C)$ có phương trình tiếp tuyến thì phương trình $ x_0^2+6x_0=k$ phải  có nghiệm $\Leftrightarrow \Delta \geq 0 \Leftrightarrow k\geq -9$.
   \item Để đồ thị  $(C)$ có $1$ phương trình tiếp tuyến duy nhất thì phương trình $ x_0^2+6x_0=k$ phải có nghiệm kép $\Leftrightarrow \Delta = 0 \Leftrightarrow k= -9$.
  \end{itemize}
  $\Leftrightarrow x_0=-3$$\Rightarrow y_0=f(x_0)=16$ 
 } 
\end{ex}

% nguồn: Ok_dungsai_lop11_Lan1
\begin{ex}%[1D7H2-4]%[ma:f42977e4]
 Cho hàm số $y=x^2+3x+1$ có đồ thị $(C)$. Viết được phương  trình tiếp tuyến của $(C)$ tại giao điểm của $(C)$ với trục tung. Khi đó
 \choiceTF{\True Hệ số góc của phương trình tiếp tuyến bằng $3$ }{Phương  trình tiếp tuyến đi qua điểm $A(1;3)$ }{\True Phương trình tiếp tuyến cắt đường thẳng $y=2x+1$ tại điểm có hoành độ bằng $0$}{\True Phương trình tiếp tuyến vuông góc với đường thẳng $y=-\dfrac{1}{3}x+1$}
 \loigiai{ 
   Với $x_0=0 \Rightarrow y_0=1$.\\
  Ta có $f'(0)=\lim \limits_{x \to 0} \dfrac{f(x)-f(0)}{x-0}=\lim\limits _{x\to 0} \dfrac{x^2+3  x}{x}=\lim\limits_{x \to 0}(x+3)=3$.\\
  Vậy phương trình tiếp tuyến là $y=3x+1$. Tiếp tuyến này không đi qua điểm $A(1;3)$ và vuông góc với  đường thẳng $y=-\dfrac{1}{3}x+1$.\\
  Xét phương trình hoành  độ giao điểm của tiếp tuyến với đường thẳng $y=2x+1$, ta có
  $$3x+1=2x+1\Leftrightarrow x=0.$$ 
  }
\end{ex}

% nguồn: Ok_dungsai_lop11_Lan1
\begin{ex}%[1D7H2-4]%[ma:ac127c86]
 Xét tính đúng, sai  của các mệnh đề sau
 \choiceTF 
  {Phương trình tiếp tuyến  của đồ thị hàm số $y=-x^3+3x^2+2$ $(C)$ mà có hệ số góc lớn nhất là $y=3x-1$}
   {\True Tiếp tuyến với đồ thị $(C)\colon y=x^3-3x^2-2$ song song với đường thẳng $(d)\colon y=9x+3$ có phương trình là $y=9x-29$.}
  {\True Số các phương trình tiếp tuyến của đồ thị hàm số $y=\dfrac{2x+1}{x-1}$ song song với đường thẳng $y=-3x+15$ là $2$} 
  {Phương trình tiếp tuyến của đồ thị hàm số $y=\dfrac{x^3}{3}+3x^2-2$ $(C)$ có tiếp tuyến  có hệ số góc $k=-9$ là $y=-9(x+3)$}
  \loigiai{
  \begin{itemchoice}
   \itemch Ta có $y'=-3x^2+6x=-3(x-1)^2+3\leqslant 3$; với  $x=1$ thì $y'=3$.\\
   Nên hệ  số góc lớn nhất của tiếp tuyến là $k=3$.\\
   Phương trình tiếp  tuyến $y=3(x-1)+4\Leftrightarrow y=3x+1$.
   \itemch $y'=3x^2-6x$.\\ 
    Hoành độ tiếp điểm là nghiệm của phương trình $3x^2-6x=9\Leftrightarrow x^2-2x-3=0\Leftrightarrow \hoac{&x=-1\\&x=3.} $\\
   Với  $x=-1$ ta có $y=-6$ phương trình tiếp tuyến là $y+6=9(x+1)\Leftrightarrow y=9x+3$.\\
   Với $x=3$ ta có $y=-2$ phương trình tiếp tuyến  là $y+2=9(x-3)\Leftrightarrow y=9x-29$.
   \itemch Gọi $M(x_0;y_0)$, $x_0\ne 1$ là  tiếp điểm
   $y'=-\dfrac{3}{(x-1)^2}$.\\ 
   Đồ thị hàm số  song song với $y=-3x+15$ nên ta có $f'(x_0)=-3\Leftrightarrow -\dfrac{3}{(x_0-1)^2}=-3 \Leftrightarrow \hoac{&x_0=0\\&x_0=2.} $
   \begin{itemize}
    \item Với $x_0=0 \Rightarrow  y_0=-1 \Rightarrow $ phương trình tiếp tuyến là $y=-3x-1$.\\
    \item Với  $x_0=2 \Rightarrow y_0=5 \Rightarrow $ phương trình tiếp tuyến là $y=-3x+11$.
   \end{itemize}
    \itemch Ta có $y'=x^2+6x$.\\
   Vì tiếp tuyến có hệ số góc  $k=-9$ nên $y'=-9$$\Leftrightarrow x^2+6x+9=0$$\Leftrightarrow x=-3$.\\
   Với $x=-3$$\Rightarrow y=16$. Vậy phương trình  tiếp tuyến cần tìm là $y-16=-9(x+3)$.
  \end{itemchoice}
  }
\end{ex}

% nguồn: Ok_dungsai_lop11_Lan1
\begin{ex}%[1D7V2-4]%[ma:925cbdd1]
 Cho hàm số $y=x^3-5x^2+2$ có  đồ thị $(C)$. Xét tính Đúng - Sai của các mệnh đề sau:
  \choiceTF
  {\True Hàm  số đã cho liên tục trên $\mathbb{R}$}
   {$y'=x^2-5x$}
  {Bất phương trình $y'<0$ có  $4$ nghiệm nguyên}
  {\True Tồn tại hai điểm $M$  trên $(C)$ thỏa mãn tiếp tuyến của $(C)$ tại điểm $M$ vuông góc với đường thẳng $y=\dfrac{1}{7}x-4$}
  \loigiai{
  \begin{itemchoice}
    \itemch Hàm số đã cho là hàm đa thức liên tục trên $\mathbb{R}$.
   \itemch $y'=3x^2-10x$. 
   \itemch $y'<0 \Leftrightarrow 3x^2-10x< 0  \Leftrightarrow x \in \left(0; \dfrac{10}{3} \right)$.\\
    Bất phương trình $y'<0$ có ba nghiệm nguyên $x \in \{1;2;3\}$.
   \itemch Vì tiếp tuyến tại $M$ vuông góc với đường thẳng $y=\dfrac{1}{7}x-4$ nên  $$(3x^2-10x)\cdot\dfrac{1}{7}=-1\Leftrightarrow 3x^2-10x+7=0\Leftrightarrow\hoac{&x=1\\&x=\tfrac{7}{3}.}$$\\
   Vậy có hai điểm $M$ thỏa  mãn yêu cầu bài toán.
 \end{itemchoice}}
\end{ex}

% nguồn: Ok_dungsai_lop11_Lan1
\begin{ex}%[1D7V2-4]%[ma:7a3a0765]
 Trong hệ trục $Oxy$, cho hàm số $y=2x^2-3x$ và đường thẳng $\Delta\colon x+y-2=0$. Gọi $d$ là tiếp tuyến của đồ thị hàm số đã cho  tại điểm có hoành độ bằng $x_0$. Các mệnh đề sau đúng hay sai?
  \choiceTF
  {Đường thẳng $\Delta$  có hệ số góc bằng $1$}
  {\True Hệ số góc của $d$  là $k=4x_0-3$}
  {Nếu $d\parallel \Delta$  thì $x_0=-1$}
  {Nếu $d\parallel \Delta$ thì phương trình của  $d$ là $d\colon 2x+2y-1=0$}
  \loigiai{
  \begin{itemchoice}
   \itemch Phương trình đường thẳng $\Delta$ được viết lại  là $y=-x+2$ và có hệ số góc bằng $-1$.
   \itemch Ta có  $y'=4x-3$ và hệ số góc của $d$ là $k=4x_0-3$.
   \itemch Vì $d\parallel \Delta$  nên $4x_0-3=-1\Leftrightarrow 4x_0=2\Leftrightarrow x_0=\dfrac{1}{2}$.
   \itemch Ta có $x_0=\dfrac{1}{2}$ nên tung độ của tiếp điểm là $y_0=2\cdot\left(\dfrac{1}{2}\right)^2-3\cdot\dfrac{1}{2}=-1$.\\  Phương trình tiếp tuyến $d$ là $y=-1\left(x-\dfrac{1}{2}\right)-1\Leftrightarrow y=-x-\dfrac{1}{2}\Leftrightarrow 2x+2y+1=0$.
  \end{itemchoice}
  }
\end{ex}

% nguồn: Ok_dungsai_lop11_Lan1
\begin{ex}%[1D7V2-4]%[ma:3279b7fe]
 Viết được phương trình tiếp tuyến cảu đồ thị hàm số $y=\dfrac{x+9}{x+1}$, biết tiếp tuyến vuông  góc với đường thẳng $d\colon x-2y+2=0$. Khi đó
  \choiceTF
  {\True Có hai phương trình tiếp tuyến  thỏa mãn}
  {\True Hệ số  góc tiếp tuyến bằng $-2$}
  {\True Phương trình  tiếp tuyến đi qua điểm $A(1;5)$}
  {Phương trình tiếp tuyến đi qua  điểm $B(1;-7)$}
 \loigiai{ 
  \begin{itemize}
    \item Ta có $y'=\dfrac{-8}{(x+1)^2}$ và đường thẳng $d\colon y=\dfrac{1}{2}x+1$.\\
   Gọi $M(x_0;y_0)$ là tiếp điểm.\\ 
   Do tiếp tuyến vuông góc với $d$ nên hệ số góc của tiếp  tuyến là
   \allowdisplaybreaks 
   \begin{eqnarray*}
    k=y'(x_0)=-2&\Leftrightarrow&\dfrac{-8}{(x_0+1)^2}=-2\\ 
    &\Leftrightarrow&(x_0+1)^2=4\\ 
    &\Leftrightarrow&\hoac{&x_0+1=2\\&x_0+1=-2}\\ 
     &\Leftrightarrow&\hoac{&x_0=1\\&x_0=-3.}
   \end{eqnarray*}
   Với  $x_0=1$, suy ra $y_0=5$. Khi đó, phương trình tiếp tuyến với đồ thị hàm số là $$y=-2(x-1)+5\Leftrightarrow y=-2x+7.$$
   Với $x_0=-3$, suy ra $y_0=-3$.  Khi đó, phương trình tiếp tuyến với đồ thị hàm số là $$y=-2(x+3)-3\Leftrightarrow y=-2x-9.$$
   \item Hệ số góc của tiếp tuyến là  $k=-2$.
   \item Phương trình tiếp tuyến qua điểm $A(1;5)$. 
   \item Phương trình tiếp tuyến không qua  điểm $A(1;-7)$.
  \end{itemize}
  }
\end{ex}

% nguồn: Ok_dungsai_lop11_Lan1
\begin{ex}%[1D7V2-4]%[ma:35e14a03]
 Cho hàm số $y=f(x)=\dfrac{3x+2}{x+1}$. Các  mệnh đề sau đúng hay sai?
 \choiceTF 
  {\True $f'(x)>0$, $\forall x\neq-1$} 
  {$f'(2)+f'(-3)=\dfrac{5}{36}$} 
  {Tích tất  cả các nghiệm của phương trình $f'(x)=25$ là $\dfrac{26}{25}$}
  {\True Tiếp tuyến của đồ thị hàm số tại điểm $M(-2;4)$ đi  qua điểm $E(-5;1)$}
  \loigiai{
  \begin{itemchoice}
   \itemch \textbf{Đúng}.\\ 
   $f'(x)=\dfrac{(3x+2)'\cdot (x+1)-(3x+2)\cdot (x+1)'}{(x+1)^2}=\dfrac{1}{(x+1)^2}>0$,  $\forall x\neq-1$.
   \itemch  \textbf{Sai}.\\
    Ta có $f'(2)+f'(-3)=\dfrac{1}{(2+1)^2}+\dfrac{1}{(-3+1)^2}=\dfrac{13}{36}$.
   \itemch  \textbf{Sai}.\\
   Ta có  $f'(x)=25\Leftrightarrow\dfrac{1}{(x+1)^2}=25\Leftrightarrow(x+1)^2=\dfrac{1}{25}\Leftrightarrow\hoac{&x_1=-\dfrac{6}{5}\\&x_2=-\dfrac{4}{5}.}$ \\
   Nên  $x_1\cdot x_2=\dfrac{24}{25}$.
    \itemch \textbf{Đúng}.\\
   Ta  có $f'(-2)=\dfrac{1}{(-2+1)^2}=1$.\\
    Phương trình tiếp tuyến của đồ thị hàm số tại điểm $M(-2;4)$ là $$y-4=1(x+2)\Leftrightarrow y=x+6.$$
   Vì  $1=-5+6$ suy ra tiếp tuyến đi qua điểm $E(-5;1)$.
 \end{itemchoice}}
\end{ex}

% nguồn: Ok_dungsai_lop11_Lan2
\begin{ex}%[1D7V2-4]%[ma:e2b055f7]
	Viết được phương trình tiếp tuyến $\Delta$ của đồ thị  hàm số $y=\dfrac{x+9}{x+1}$ biết tiếp tuyến vuông góc với đường thẳng $d\colon x-2y+2=0$.
	\choiceTF 
		{\True Có hai phương trình tiếp tuyến  thỏa mãn}
		{\True Hệ số góc của tiếp tuyến bằng  $-2$}
		{\True Phương trình tiếp tuyến $\Delta$ đi qua  điểm $A(1;5)$}
		{\True Phương trình tiếp tuyến $\Delta$ đi  qua điểm $B(1;-7)$}
	\loigiai{ 
		Đường thẳng $d\colon x-2y+2=0\Leftrightarrow y=\dfrac{1}{2} x+1$ nên  đường thẳng $d$ có hệ số góc là $k_d=\dfrac{1}{2}$.
		\begin{itemchoice}
			\itemch Vì tiếp tuyến $\Delta$ vuông góc  với đường thẳng $d$ nên
			\[k_{\Delta}\cdot  k_d=-1\Rightarrow k_{\Delta}=-\dfrac{1}{k_d}=-2.\]
			Hoành độ tiếp điểm là nghiệm của phương trình $y'=k_{\Delta}\Rightarrow  \dfrac{-8}{(x+1)^2}=-2\Rightarrow \hoac{&x=1\\&x=-3.}$
			\begin{itemize}
				\item Với $x=1$ ta  có phương trình tiếp tuyến $\Delta_1$ là $y=-2x+7$.
				\item Với $x=-3$ ta  có phương trình tiếp tuyến $\Delta_2$ là $y=-2x-9$.
			\end{itemize}
			Vậy có hai phương trình tiếp tuyến thỏa mãn là $\Delta_1 \colon y=-2x+7$;  $\Delta_2\colon y=-2x-9$.
			\itemch Vì tiếp  tuyến $\Delta$ vuông góc với đường thẳng $d$ nên
			 \[k_{\Delta}\cdot k_d=-1\Rightarrow k_{\Delta}=-\dfrac{1}{k_d}=-2.\]
			\itemch Vì $5=-2\cdot 1+7$  nên $A\in\Delta_1$. \\
			Vậy tiếp tuyến $\Delta$ đi qua điểm  $A(1;5)$.
			\itemch Vì $\heva{&-7=-2\cdot 1+7\text{ sai}\\&-7=-2\cdot 1-9\text{ sai}}$ nên $B\notin\Delta_1$ và $B\notin\Delta_2$. \\ 
			Vậy các tiếp tuyến $\Delta_1$  và $\Delta_2$ không đi qua điểm $B$.
		\end{itemchoice}
	} 
\end{ex}

% nguồn: Ok_dungsai_lop11_Lan2
\begin{ex}%[1D7V2-4]%[ma:9deda153]
	 Cho hàm số $y=f(x)=x^3-3x^2-9x+10$ có đồ thị $(C)$.
	\choiceTF 
		{$f'(x)=3x^2-3x-9$} 
		{\True Tập nghiệm của bất phương  trình $f'(x) < 0$ là $S=(-1;3)$}
		 {\True Phương trình $f'(x)=0$ có tổng các nghiệm bằng $2$}
		{Tiếp tuyến của $(C)$ có hệ số góc  nhỏ nhất có phương trình là $y=12x-1$}
	\loigiai{ 
		\begin{itemchoice}
			\itemch $f'(x)=3x^2-6x-9$. 
			\itemch $f'(x)=3x^2-6x-9< 0\Leftrightarrow-1<  x < 3$.
			 \itemch $f'(x)=0\Leftrightarrow 3x^2-6x-9=0\Leftrightarrow \hoac{&x=-1\\&x=3.}$
			\itemch Gọi $x_0$ là tiếp điểm, hệ số góc của tiếp tuyến tại  $x_0$ là
			 \[f'(x_0)=3x_0^2-6x_0-9=3\left(x_0-1\right)^2-12\ge-12
			\] 
			Hệ số góc nhỏ nhất của tiếp tuyến bằng $-12$  tại $x_0=1\Rightarrow y_0=-1$.\\
			Vậy tiếp tuyến của $(C)$ có hệ số góc nhỏ nhất có phương trình \[y=-12(x-1)-1\Leftrightarrow y=-12x+11.\] 
		\end{itemchoice}
	} 
\end{ex}

% nguồn: Ok_dungsai_lop11_Lan2
\begin{ex}%[1D7V2-4]%[ma:c246c731]
	Cho hàm số $f(x)=\dfrac{1}{2}{x^4}-3x^2+2\,024$ có đồ thị $(C)$ và đường thẳng  $d\colon y=-4x+\dfrac{4043}{2}$.
	 \choiceTF
		{\True Đạo hàm của $f(x)$  là $f'(x) = 2x^3-6x$}
		{Hệ số góc của đường thẳng $d$ là  $4$}
		{\True Gọi $\Delta $ là tiếp tuyến của đồ thị $(C)$ song song với $d$, khi đó hệ số góc của $\Delta  $ là $-4$}
		{Có  hai tiếp tuyến của đồ thị $(C)$ song song với $d$}
	 \loigiai{
		\begin{itemchoice}
			\itemch Ta có  $ f'(x)=2x^3-6x $.
			\itemch Hệ số góc của đường thẳng $d$  là $-4$.
			\itemch Do $ \Delta \parallel d $ nên hệ số góc của đường thẳng  $ \Delta $ bằng hệ số góc của đường thẳng $ d $ là $ -4 $.
			\itemch  Ta có $f'(x)=2x^3-6x$.\\
			Gọi  $M\left(x_0,y_0\right)$ là tọa độ tiếp điểm.\\
			Do tiếp  tuyến $\Delta $ song song với đường thẳng $d$ nên $f'\left(x_0\right)=-4\Leftrightarrow 2x_0^3-6x_0=-4 \Leftrightarrow \hoac{&x_0=-2\\&x_0=1.}$\\
			Với $x_0=1\Rightarrow{y_0}=\dfrac{4043}{2}$. Phương trình tiếp tuyến $y=-4x+\dfrac{4043}{2}$  (loại) vì trùng $d$.\\
			Với $x_0=-2\Rightarrow{y_0}=2\,020$.\\ 
			Phương trình  tiếp tuyến $y=-4x+2012$ (nhận).
		\end{itemchoice}
	 }
\end{ex}


%%%%%%%%%%%%%%%% TRẢ LỜI NGẮN %%%%%%%%%%%%%%%%

% nguồn: Ok_TraLoiNgan_Lop11_Lan2
\begin{ex}%[1D7H2-4]%[ma:6e5b7e38]
	Cho biết  điện lượng truyền trong dây dẫn theo thời gian biểu thị bởi hàm số $Q(t)=2t^2+t$, trong đó $t$ được tính bằng giây và $Q$ được tính theo Cu-lông. Tính cường độ dòng điện tại thời điểm $t=4$\,(s).
	\shortans{17} 
	\loigiai{ 
		Ta có $Q' (t)=\left(2t^2+t\right)'=4t+1$.\\ 
		Do đó cường độ dòng điện tại thời điểm $t=4$\,(s) là $Q'  (4)=4\cdot 4+1=17$\,(A).
	 }
\end{ex}

% nguồn: Ok_TraLoiNgan_Lop11_Lan2
\begin{ex}%[1D7H2-4]%[ma:6eb93fb2]
	Cho đường cong $(C)\colon y=\dfrac{2x-1}{x-1}$. Có bao nhiêu phương trình tiếp tuyến với $(C)$ biết hệ số góc của tiếp tuyến  bằng $-\dfrac{1}{4}$.
	\shortans{2} 
	\loigiai{ 
		Ta có $y'=-\dfrac{1}{(x-1)^2}$ và  $y(2)=3$.\\
		Hệ số góc của tiếp tuyến  bằng $-\dfrac{1}{4}$.\\
		 Nên $-\dfrac{1}{\left(x_0-1\right)^2}=-\dfrac{1}{4}\Rightarrow{\left(x_0-1\right)^2}=4\Rightarrow \hoac{&x_0=3\\&x_0=-1.}$
		\begin{itemize}
			\item Với $x_0=3\Rightarrow{y_0}=\dfrac{5}{2}$. Phương trình tiếp tuyến  là $$y=-\dfrac{1}{4}(x-3)+\dfrac{5}{2}=-\dfrac{1}{4}x+\dfrac{13}{4}.$$
			\item Với $x_0=-1\Rightarrow{y_0}=\dfrac{3}{2}$ . Phương trình tiếp tuyến là 
			 $$y=-\dfrac{1}{4}(x+1)+\dfrac{3}{2}=-\dfrac{1}{4}x+\dfrac{5}{4}.$$
	\end{itemize}}
\end{ex}

% nguồn: Ok_TraLoiNgan_Lop11_Lan1
\begin{ex}%[1D7V2-4]%[ma:6d4dd744]
	Đồ thị hàm số $y=\dfrac{ax+b}{x-1}$  cắt trục tung tại điểm $A(0;-1)$, tiếp tuyến của đồ thị tại điểm $A$ có hệ số góc $k=-3$. Giá trị của $a+b$ là
	\shortans{3} 
	 \loigiai{
		Đồ thị hàm số $y=\dfrac{ax+b}{x-1}$ cắt trục tung tại điểm $A(0;-1)$ nên $\dfrac{a\cdot  0+b}{0-1}=-1\Leftrightarrow b=1$.\\
		Khi đó $y=\dfrac{ax+1}{x-1}\Rightarrow  y'=\dfrac{-a-1}{(x-1)^2}\Rightarrow y'(0)=-a-1$.\\
		Do tiếp tuyến của đồ thị tại điểm $A$ có hệ số  góc $k=-3$ nên $-a-1=-3\Leftrightarrow a=2$.\\
		 Suy ra $a=2; b=1$. Vậy $a+b=3$.
	} 
\end{ex}

% nguồn: Ok_TraLoiNgan_Lop11_Lan1
\begin{ex}%[1D7V2-4]%[ma:e9c4fb84]
	Gọi $\Delta$ là tiếp tuyến của đồ  thị hàm số $y = -x^3+3x^2-3x-2$ song song với đường thẳng $y = -3x-2$. Tìm tung độ giao điểm của $\Delta$ với trục tung.
	\shortans{2} 
	\loigiai{ 
	Ta có  $y' = -3x^2 + 6x - 3$.\\
	Gọi $x_0$ là hoành độ tiếp điểm.\\ 
	Vì tiếp tuyến của đồ thị hàm số song song với đường thẳng $y = -3x-2$ nên 
	\begin{align*}
		y'(x_0) = -3 \Leftrightarrow -3x_0^2 + 6x_0 - 3 = - 3 \Leftrightarrow -3x_0^2 + 6x_0 = 0 \Leftrightarrow \left[\begin{aligned}& x_0 = 0 \\& x_0 = 2.
	\end{aligned}\right.
\end{align*}
\begin{itemize}
	\item Với $x_0 = 0$ thì $y_0 = -2$. Phương trình tiếp tuyến của đồ thị là $y = -3x-2$ (không thỏa mãn vì tiếp tuyến này  trùng với đường thẳng đã cho).
	\item Với $x_0 = 2$  thì $y_0 = -4$. Phương trình tiếp tuyến của đồ thị là $y = -3x+2$.
\end{itemize}
Vậy $\Delta\colon y = -3x+2$. Do đó tung độ giao điểm của $\Delta$ với trục tung  bằng $2$.
 }
\end{ex}

% nguồn: Ok_TraLoiNgan_Lop11_Lan1
\begin{ex}%[1D7V2-4]%[ma:a47db10c]
	Cho hàm số $f(x)=x^3+3x+1$ có đồ thị $(C)$ và đường  thẳng $d\colon 6x-y-1=0$. Biết rằng phương trình tiếp tuyến của $(C)$ song song với đường thẳng $d$ có dạng $y=6x+b$. Tìm giá trị của $b$.
	\shortans{3} 
	\loigiai{ 
		Ta có:  $y'=3x^2+3$.\\
		Do  tiếp tuyến song song với đường thẳng $6x-y-1=0$ nên $y'(x_0)=6\Leftrightarrow x_0=\pm 1$.\\
		 Với $x_0=1$ thì $y_0=5$. Suy ra $y=6x-1$ (loại).\\
		 Với $x_0=-1$ thì $y_0=-3$. Suy ra $y=6x+3$ (thỏa).\\
		Vậy phương trình tiếp tuyến cần tìm  là $y=6x+3$, do đó $b=3$.
	} 
\end{ex}

% nguồn: Ok_TraLoiNgan_Lop11_Lan2
\begin{ex}%[1D7V2-4]%[ma:a6ae7d35]
	Gọi $M\left(x_0;y_0\right)$ là điểm trên đồ thị hàm số  $y=x^3-3x^2-1$ mà tiếp tuyến tại đó có hệ số góc bé nhất trong các tiếp tuyến của đồ thị hàm số. Khi đó $x_0^2+y_0^2$ bằng bao nhiêu?
	\shortans{10} 
	 \loigiai{
		Ta  có $y'=3x^2-6x$.\\
		Suy ra hệ số góc của tiếp  tuyến tại $M\left(x_0;{y_0}\right)$ là $f'\left(x_0\right)=3x_0^2-6x_0$.\\
		 Ta có $f'\left(x_0\right)=3x_0^2-6x_0+3-3=3\left(x_0-1\right)^2-3\ge-3$.\\
		Suy ra hệ  số góc của tiếp tuyến tại $M\left(x_0;{y_0}\right)$ nhỏ nhất bằng $-3$ khi
		$$  x_0-1=0 \Leftrightarrow x_0=1 \Rightarrow y_0=-3. $$
	 Khi đó $x_0^2+y_0^2=1^2+\left(-3\right)^2=10$.}
\end{ex}

% nguồn: Ok_TraLoiNgan_Lop11_Lan2
\begin{ex}%[1D7V2-4]%[ma:d4794bcf]
	Cho hàm số $y=x^3-3x^2+1$ có đồ thị $(C)$. Phương trình tiếp tuyến của đồ  thị $(C)$ song song với đường thẳng $\Delta\colon 9x-y+6=0$ có dạng $y=ax+b$. Tính $5a+3b$.
	\shortans{27} 
	\loigiai{ 
		Đường thẳng $\Delta\colon  9x-y+6=0\Rightarrow y=9x+6$ có hệ số góc là $9$.\\
		Vì tiếp tuyến cần tìm song song với đường thẳng $\Delta$  nên tiếp tuyến có hệ số góc $k=9$.\\
		Suy ra hoành độ tiếp điểm là nghiệm của phương trình  $y'=k\Rightarrow 3x^2-6x=9\Rightarrow \hoac{&x=3\\&x=-1.}$
		\begin{itemize}
			\item Với $x=-1$, phương trình tiếp tuyến là $y=9x+6$  (loại vì trùng với đường thẳng $\Delta$).
			\item Với $x=3$, phương trình tiếp tuyến là  $y=9x-26$ (thỏa mãn).
		\end{itemize}
		Vậy phương trình  tiếp tuyến cần tìm là $y=9x-6$.\\
		 Suy ra $a=9$; $b=-6$. Do đó $5a+3b=27$.
	} 
\end{ex}

% nguồn: Ok_TraLoiNgan_Lop11_Lan2
\begin{ex}%[1D7V2-4]%[ma:6943a87e]
	Cho hàm số $y=-x^3+3x^2+9x-1$ có đồ thị là $(C)$. Tìm hệ số góc lớn nhất của  tiếp tuyến tại một điểm $M$ trên đồ thị $(C)$.
	\shortans{12} 
	\loigiai{ 
		Gọi điểm $M\left(x_0;{y_0}\right)\in (C)$ là toạ  độ tiếp điểm và $y'=-3x^2+6x+9$.\\
		Hệ số góc của tiếp  tuyến tại điểm $M$ là $y'\left(x_0\right)=-3x_0^2+6x_0+9$.\\
		Ta thấy, hệ số góc của tiếp tuyến tại  điểm $M$ là một hàm số có đồ thị là một parabol, có đỉnh $S\left(-\dfrac{b}{2a};-\dfrac{\Delta}{4a}\right)\Rightarrow S(1;12)$ và hệ số $a=-3 < 0$ nên hàm số có giá trị lớn nhất bằng $ 12 $ tại $x_0=1$.\\
	Vậy hệ số góc  lớn nhất của tiếp tuyến là $y'(1)=12$.}
\end{ex}


%%%%%%%%%%%%%%%% TỰ LUẬN %%%%%%%%%%%%%%%%

% nguồn: Ok_TuanLuan_Toan11_Lan1
\begin{bt}%[1D7H2-4]%[ma:7b114570]
	Cho hàm số $y=\dfrac{1}{3} x^{3}+\dfrac{1}{2} x^{2}-2 x+1$ (đồ thị là $(C)$). Viết phương trình tiếp tuyến  của $(C)$ biết hệ số góc của tiếp tuyến bằng $4$.
	 \loigiai{
		\begin{itemize}
			\item  Đạo hàm của hàm số $y=\dfrac{1}{3} x^{3}+\dfrac{1}{2} x^{2}-2x+1$ là $y'=x^2+x-2$.
			\item Gọi $M\left(x_0;y_0\right)$ là tọa độ tiếp điểm. \\ 
			Vì tiếp tuyến có hệ số góc bằng  $4$ nên ta có phương trình\\
			\begin{align*}
				 &y'\left(x_0\right)=4\\
				 \Leftrightarrow& x_0^2+x_0-2=4\\
				\Leftrightarrow  & x_0^2+x_0-6=0\\
				 \Leftrightarrow & \hoac{&x_0=2\\&x_0=-3.}
			\end{align*}
			 \item Với $x_0=2 \Rightarrow y_0= \dfrac{5}{3}$.\\
			Ta được  phương trình tiếp tuyến $y=4(x-2)+\dfrac{5}{3} \Leftrightarrow 12x-3y -19=0$.
			\item Với $x_0= -3 \Rightarrow y_0=  \dfrac{5}{2}$.\\
			Ta được  phương trình tiếp tuyến $y=4(x+3)+\dfrac{5}{2} \Leftrightarrow 8x-2y + 29=0$.
		\end{itemize}
		Vậy có hai phương trình tiếp tuyến thỏa  mãn yêu cầu bài toán là
		\[12x-3y -19=0\, \,\text{và } \,\,8x-2y + 29=0.\] 
	} 
\end{bt}

% nguồn: Ok_TuanLuan_Toan11_Lan2
\begin{ex}%[1D7H2-4]%[ma:b1300a12]
	Cho hàm số $y=\dfrac{1}{3} x^{3}+\dfrac{1}{2} x^{2}-2 x+1$ (đồ thị là $(C)$). Viết phương trình tiếp tuyến của $(C)$  biết hệ số góc của tiếp tuyến bằng $4$.
	\loigiai{ 
		\begin{itemize}
			\item Đạo hàm của hàm số  $y=\dfrac{1}{3} x^{3}+\dfrac{1}{2} x^{2}-2x+1$ là $y'=x^2+x-2$.
			 \item Gọi $M\left(x_0;y_0\right)$ là tọa độ tiếp điểm. \\
			Vì tiếp tuyến có hệ số góc  bằng $4$ nên ta có phương trình\\
			\begin{align*}
				 &y'\left(x_0\right)=4\\
				\Leftrightarrow& x_0^2+x_0-2=4\\ 
				 \Leftrightarrow & x_0^2+x_0-6=0\\
				\Leftrightarrow  & \hoac{&x_0=2\\&x_0=-3.}
			\end{align*}
			 \item Với $x_0=2 \Rightarrow y_0= \dfrac{5}{3}$.\\
			Ta  được phương trình tiếp tuyến $y=4(x-2)+\dfrac{5}{3} \Leftrightarrow 12x-3y -19=0$.
			\item Với $x_0= -3 \Rightarrow y_0= \dfrac{5}{2}$.\\ 
			Ta được phương trình tiếp tuyến $y=4(x+3)+\dfrac{5}{2}  \Leftrightarrow 8x-2y + 29=0$.
		\end{itemize}
		Vậy có hai phương  trình tiếp tuyến thỏa mãn yêu cầu bài toán là
		 \[12x-3y -19=0\, \,\text{và } \,\,8x-2y + 29=0.\]
	 }
\end{ex}

% nguồn: Ok_TuanLuan_Toan11_Lan2
\begin{ex}%[1D7H2-4]%[ma:8f6683fb]
	Cho đường cong $(C)\colon y=x^3-3x^2-2x-1$ và đường thẳng $d\colon 2x+y+1=0$.  Viết phương trình tiếp tuyến của đường cong $(C)$ biết tiếp tuyến song song với $d$.
	\loigiai{ 
		Ta  có đường thẳng $2x+y+1=0\Leftrightarrow y=-2x-1$.\\
		Tiếp tuyến song song với $d$ nên tiếp tuyến có hệ số góc là  $k=-2$.\\
		Từ $y=x^3-3x^2-2x-1$ suy ra  $y'=3x^2-6x-2$.\\
		Gọi $x_0$ là  hoành độ của tiếp điểm.
		Khi đó $$y'(x_0)=k\Leftrightarrow 3x_0^2-6x_0-2=-2\Leftrightarrow  3x_0^2-6x_0=0\Leftrightarrow\hoac{&x_0=0\\&x_0=2.}$$
		Với $x_0=0$  thì $M(0;-1)$, phương trình tiếp tuyến là $y=-2x-1$ trùng với $d$.\\
		Với $x_0=2$ thì $N(2;-9)$, phương  trình tiếp tuyến là $y=-2x-5$.\\
		Vậy phương trình  tiếp tuyến cần tìm là $y=-2x-5$.
	 }
\end{ex}

% nguồn: Ok_TuanLuan_Toan11_Lan2
\begin{ex}%[1D7H2-4]%[ma:360f64ba]
	 Cho hàm số $y=x^2-2x+2023$ có đồ thị $(C)$ và đường thẳng $d\colon y=2x+7$. Viết phương trình tiếp tuyến của $(C)$ biết tiếp tuyến song song với $d\colon y=2x+7$.
	 \loigiai{
		Gọi điểm $M(x_0; y_0)$ là tiếp điểm giữa  tiếp tuyến và đồ thị $(C)$.\\
		Ta có $y'=2x  - 2$ suy ra hệ số góc $k=y'(x_0)=2x_0-2$.\\
		Mặt khác tiếp tuyến song  song với đường thẳng $d\colon y=2x+7$ nên $$k=2\Leftrightarrow 2x_0-2=2\Leftrightarrow x_0=2\Rightarrow y_0=2023.$$
		Như thế, $M(2;2023)$.\\ 
		Vậy phương trình tiếp  tuyến là $y=2(x-2)+2023=2x+2019$.
	 }
\end{ex}

% nguồn: Ok_TuanLuan_Toan11_Lan2
\begin{ex}%[1D7H2-4]%[ma:159f71f8]
	 Cho hàm số $y=x^4+2x^2-1$ có đồ thị $(C)$. Viết phương trình tiếp tuyến của $(C)$ biết tiếp tuyến vuông góc với đường thẳng $d\colon y=\dfrac{x}{8}-2022$.
	\loigiai{ 
		Ta có $y'=4x^3+4x$. Suy ra hệ số góc của tiếp tuyến là  $k=4x^3+4x$.\\
		Vì tiếp tuyến vuông góc với đường thẳng $d\colon y=\dfrac{x}{8}-2022$ nên $k\cdot\dfrac{1}{8}=-1\Leftrightarrow  k=-8$.\\
		Hoành độ tiếp điểm là nghiệm của phương trình $$4x^3+4x=-8\Leftrightarrow x^3+x-2=0\Leftrightarrow  x=1.$$
		Với $x=1$ thì tọa độ tiếp điểm  $M(1;2)$.\\
		Vậy phương trình tiếp tuyến là $y=8(x-1)+2$  hay $y=8x-6$.
	} 
\end{ex}

% nguồn: Ok_TuanLuan_Toan11_Lan2
\begin{ex}%[1D7H2-4]%[ma:b8ec914c]
	Cho hàm số $y=2x^3-3x^2-1$ có đồ thị $(C)$. Viết phương trình tiếp tuyến  của $(C)$ biết tiếp tuyến song song với đường thẳng $d\colon y=12x+2022$.
	 \loigiai{Ta có $y'=6x^2-6x$.\\
		Hoành độ tiếp điểm là  nghiệm của phương trình $$6x^2-6x=12\Leftrightarrow x^2-x-2=0\Leftrightarrow\hoac{&x=-1\\&x=2.}$$
		Với $x=-1$ thì tọa độ tiếp điểm $M(-1;-6)$. Do đó phương trình tiếp tuyến $d_1\colon y=12x+6$.\\ 
		Với $x=2$ thì tọa độ  tiếp điểm $N(2;3)$. Suy ra phương trình tiếp tuyến là $d_2\colon y=12x-21$.\\
		Dễ thấy $d_1\parallel d$ và  $d_2\parallel d$ nên $d_1$, $d_2$ là hai tiếp tuyến cần tìm.
	} 
\end{ex}

% nguồn: Ok_TuanLuan_Toan11_Lan2
\begin{ex}%[1D7H2-4]%[ma:97dbfca1]
	Cho hàm số $y=x^3-3x^2+2$  có đồ thị $(C)$. Viết phương trình tiếp tuyến của đồ thị hàm số $(C)$ biết tiếp tuyến song song với đường thẳng $y=9x+7$.
	 \loigiai{
		Ta có $y'=3x^2-6x$ nên hệ số  góc của tiếp tuyến là $k=3x^2-6x$.\\
		Vì tiếp tuyến song  song với đường thẳng $y=9x+7$ nên
		$$k=9\Leftrightarrow 3x^2-6x=9\Leftrightarrow 3x^2-6x-9=0\Leftrightarrow\hoac{&x=-1\\&x=3.}$$ 
		Với $x_0=-1$ thì $y_0=-2$. Phương trình tiếp tuyến là $y=9(x+1)-2\Leftrightarrow y=9x+7$ (loại vì trùng với đường  thẳng đề bài cho).\\
		Với $x_0=3$ thì $y_0=2$. Phương trình tiếp tuyến là $y=9(x-3)+2\Leftrightarrow y=9x-25$. Tiếp tuyến này song song với đường  thẳng $y=9x+7$.
	 }
\end{ex}
